\section{The continuum existence problem}\label{sec:introduction}
A phase transition in a continuum particle system can be expressed as a loss of
regularity of its thermodynamic free energy. We consider an existence problem:
can a single species of particles in $\R^3$, interacting through a stable radial
pair potential, have a temperature singularity at one fixed density, and can the same
critical inverse temperature work throughout an open interval of densities?
We use the following precise class of potentials and normalization.

A potential $\Phi(x)=\phi(|x|)$ is \emph{admissible} if it is Borel measurable,
$\Phi(0)=+\infty$, and $\phi$ is finite and locally bounded on $(0,\infty)$,
with the following properties:
\begin{enumerate}[label=(\roman*),leftmargin=*]
\item $\phi(r)\to+\infty$ as $r\downarrow0$;
\item there are $0<r_1<r_2$ and $a>0$ such that
      $\phi(r)\le-a$ for every $r\in[r_1,r_2]$;
\item $\int_R^\infty r^2|\phi(r)|\dd r<\infty$ for some finite $R$;
\item there is $B<\infty$ such that, for every finite configuration,
\begin{equation}\label{eq:stability-target}
 U_N^\Phi(x_1,\ldots,x_N):=\sum_{1\le i<j\le N}\Phi(x_i-x_j)\ge-BN.
\end{equation}
\end{enumerate}
Let $\Lambda_L=[-L/2,L/2]^3$. With free boundary conditions, define
\begin{equation}\label{eq:canonical-definition}
 Z_{\Lambda,N}^{\Phi}(\beta)=\frac1{N!}
 \int_{\Lambda^N}e^{-\beta U_N^\Phi(x_1,\ldots,x_N)}
 \dd x_1\cdots\dd x_N,\qquad \beta>0,
\end{equation}
where $Z_{\Lambda,0}^{\Phi}=1$ and $e^{-\beta\infty}=0$. The position
measure is Lebesgue measure, independent of $\beta$; no temperature-dependent
normalization is included.
For a density $\rho>0$, the canonical free-energy density is
\begin{equation}\label{eq:free-energy-definition}
 f_\Phi(\beta,\rho)=-\frac1\beta
 \lim_{L\to\infty}L^{-3}\log
 Z_{\Lambda_L,\lfloor\rho L^3\rfloor}^{\Phi}(\beta),
\end{equation}
when the limit exists. The problem is to find an admissible $\Phi$ and a fixed
$\rho$ for which this limit is finite on an interval of positive inverse
temperatures and fails to be $C^1$ across an interior point.

Write $D_\beta^-$ and $D_\beta^+$ for the left and right derivatives
with respect to inverse temperature.

\begin{theorem}\label{thm:main}
There exist an admissible radial difference-only pair potential
$\Phi(x)=\phi(|x|)$ on $\R^3$, a nonempty open interval
$I\subset(0,\infty)$, and one $\beta_c\in(1/2,3/2)$ with the following
properties. The potential satisfies $\phi(r)=o(r^{-3})$ as $r\to\infty$.
The free-cube limit \eqref{eq:free-energy-definition} exists and is finite
for every $\beta>0$ and every density $\rho>0$. For every $\rho\in I$,
the finite one-sided derivatives at the same critical point satisfy
\[
 D_\beta^- f_\Phi(\beta_c,\rho)
 >D_\beta^+ f_\Phi(\beta_c,\rho).
\]
Thus one potential and one critical inverse temperature give a strict
canonical temperature singularity throughout $I$.
\end{theorem}

The construction uses a narrow low-energy shell inside a repulsive core,
together with a summable family of radial attractions at successively larger
ranges. Thus the theorem concerns an engineered potential in the explicit
integrable-tail class above. It does not assert a transition for the ordinary
$12$--$6$ Lennard--Jones potential, and it makes no identification of the
coexisting states as conventional liquid, vapor, or solid phases.

\subsection{Context and the role of the fixed density}
The continuum phase-transition existence question belongs to the program
highlighted by Simon's problems in mathematical physics~\cite{Simon1984}.
Here the tail hypothesis is exactly condition~(iii). The construction also
gives $\phi(r)=o(r^{-3})$, but does not establish a bound
$|\phi(r)|\le C r^{-3-\varepsilon}$ with any fixed $\varepsilon>0$.
Accordingly, the theorem does not claim to resolve formulations of the
historical problem that require such a power-law margin.

Earlier rigorous continuum transitions include Ruelle's result for the
two-species Widom--Rowlinson model~\cite{Ruelle1971}. Israel's convex
construction also applies to hard-core continuum particles: a
spherically symmetric pair perturbation and a chemical potential produce
two invariant equilibrium states of different densities
\cite[Section~6, Theorem~8(a)]{Israel1975}.
The hard core in that result excludes a positive interval of separations;
the potential constructed here is finite at every positive separation.
Lebowitz, Mazel, and Presutti later proved a transition for a model with
attractive pair and repulsive four-body interactions
\cite{LebowitzMazelPresutti1998,LebowitzMazelPresutti1999}.
The four-body term remains part of that model.

Recent constructions impose other structures on the interaction.
He, Jauslin, Lebowitz, and Peled prove distinct-density Gibbs coexistence
for suitable parameters and sufficiently large range in a model with a
box-modified Kac pair interaction
\cite[Equation~(1.13) and Theorem~1.3]{HeJauslinLebowitzPeled2025}.
It depends on a fixed partition into cubes, and the reference interaction
acts only within each cube. Dereudre and Renaud-Chan prove
low-temperature coexistence for saturated finite-range interactions
satisfying a Peierls condition~\cite[Theorem~1]{DereudreRenaudChan2026}.
Their radial-kernel example integrates over unions of particle
neighborhoods, giving a many-body energy rather than an ordinary pair sum
\cite[Section~2.3, Example~3, Equation~(2.18)]{DereudreRenaudChan2026}.
Our final interaction is a radial function of the particle difference,
with all cross-cube interactions retained. Moreover, density coexistence
alone does not imply a temperature corner at a fixed density. We obtain
two supporting slopes with the same density coordinate and different
energy coordinates, the information needed for the strict derivative
jump in~\eqref{eq:free-energy-definition}.

The long-range comparison is a version of the Lebowitz--Penrose
mean-field limit~\cite{LebowitzPenrose1966}; its extension to superstable
interactions is classical~\cite{GerardiEtAl1973}. Our estimates include
local occupation and tail-continuity bounds uniform in the previously
chosen attraction ranges. The Kac limit itself fixes each finite reference
before sending the new range to infinity. The local quadratic energy
estimates belong to Ruelle's superstability framework~\cite{Ruelle1970}.
Convex methods for constructing coexistence have antecedents in
Israel's tangent-functional argument~\cite[Theorem~1]{Israel1975}
and Ruelle's restoration of lattice-gas coexistence by small changes of
long-range pair couplings and chemical potential
\cite[Section~5]{Ruelle1977Coexistence}.
Here the scaled perturbation argument keeps three separated groups of
joint pressure slopes, and hence a common ball of slopes, through
infinitely many physical approximations.

A companion construction~\cite[Theorem~1.1]{OpenAIAlgebraic2026}
gives a bounded continuous radial potential with
$|\phi(r)|\le Cr^{-3-1/32}$ for $r\ge1$ and a common temperature
singularity throughout a density interval.
That result uses a different potential and a quantitative range estimate;
the present proof retains the divergent core and proves its own
integrable-tail conclusion.


\subsection{How the construction works}
For orientation, fix a scale $t>0$ and write the activity as $e^{th}$.
The normalized pressure is the thermodynamic limit of
\[
 \frac1{tL^3}\log\sum_{N\ge0}e^{thN}Z_{\Lambda_L,N}^\Phi(\beta).
\]
It is a convex function of $(\beta,h)$. A supporting vector $(s,d)$ at
$(\beta_0,h_0)$ gives an affine plane of slope $(s,d)$ lying below that
function and touching it there. The coordinate $d$ represents density.
The proof will preserve a two-dimensional ball of these vectors, so that
one density is compatible with more than one first-coordinate slope.

The proof follows four steps.
\begin{enumerate}[leftmargin=*]
\item \textbf{Build a stable reference gas.}
The core allows pairs at an extremely small separation shell at no energy
cost. At most four particles can have all their pair distances in that shell
in $\R^3$. This geometric fact supplies a quadratic local lower bound for the
energy. Section~\ref{sec:model} uses that bound to establish the canonical
thermodynamic limit and its dual grand-canonical pressure. The scaled
function $\beta f_\Phi(\beta,\rho)/t$ is concave in inverse temperature.
\item \textbf{Create three competing pressure branches.}
For a large scale parameter $t$, the grand-canonical activity is $e^{th}$.
The reference pressure has three limiting branches, corresponding to empty
space, well-separated particles, and tetrahedral clusters. Adding a spread-out radial attraction contributes a positive quadratic
density term to the variational formula for pressure. Sections~\ref{sec:kac}
and~\ref{sec:triple} establish this comparison and produce a pressure with
three noncollinear groups of supporting slopes.
\item \textbf{Preserve those slopes at finite ranges.}
Uniform approximation alone can smooth a corner. Section~\ref{sec:convex}
therefore works at the scale of each new attractive mass: the physical
approximation error is smaller than that scale, while the density quadratic
separates the three groups. At each stage this yields a nearby point with the
same ball of supporting slopes.
\item \textbf{Pass to one potential and a density interval.}
Section~\ref{sec:construction} chooses the attraction masses before the
corresponding approximation ranges. Their summability gives one admissible
potential, and summable shifts give one limiting pressure point. The common
ball of supporting slopes passes to that physical pressure. Every horizontal slice near the center of the ball fixes a positive
density and gives two different inverse-temperature slopes of the canonical
free energy at the same critical point.
\end{enumerate}

The elementary convex facts used in these steps are proved in
Appendix~\ref{app:convex-facts}. They include local uniform convergence,
partial conjugacy, and representation of a subdifferential by limits of
ordinary gradients.
